\section{Trace monotonicity of the polar phase}\label{sec:phase}

This section is independent of the rest of the paper. Let $\Omega$ be a
complex $n\times n$ matrix with positive definite imaginary part
$\Ima\Omega=(\Omega-\Omega^\dagger)/2\ii$, where ${}^\dagger$ is the conjugate
transpose. Then $\Omega$ is invertible, and the unitary factor $\mathsf W$ in
the polar decompositions
$\Omega=(\Omega\Omega^\dagger)^{1/2}\mathsf W=\mathsf W(\Omega^\dagger\Omega)^{1/2}$
has its spectrum in the open upper half of the unit circle. The \emph{polar
phase} of $\Omega$ is the Hermitian matrix $\Phi$ with spectrum in $(0,\pi)$
and $\mathsf W=e^{\ii\Phi}$, and $|\Omega|=(\Omega\Omega^\dagger)^{1/2}=\Omega e^{-\ii\Phi}$
is the left modulus. The set of matrices with a prescribed positive definite
imaginary part is an affine space over the Hermitian matrices, and on it the
polar phase is a real-analytic map: the unitary polar factor of an invertible
matrix depends analytically on it, and $\Phi$ is $-\ii$ times its principal
logarithm.

\begin{theorem}[Monotonicity of the polar phase]\label{thm:phase}
Let $\Omega',\Omega''$ be matrices with the same positive definite imaginary
part, and let $\Phi',\Phi''$ be their polar phases. Then
\[
 \tr\bigl((\Omega''-\Omega')(\Phi''-\Phi')\bigr)\leq0,
\]
with equality only if $\Omega'=\Omega''$. More precisely, for every smooth
family $\Omega(\varepsilon)$ with constant positive definite imaginary part,
$\tr(\dot\Omega\,\dot\Phi)<0$ wherever $\dot\Omega\neq0$.
\end{theorem}

In other words, $\Rea\Omega\mapsto-\Phi$ is a strictly monotone operator
\cite{min62} on the real Hilbert space of Hermitian matrices with the trace
inner product.
For $n=1$ this says that $\arg(x+\ii y)$ decreases in $x$ when $y>0$, and for
a scalar imaginary part it follows from that by the functional calculus. In
general $\Omega$ is not normal and $\Phi$ is not a function of its
eigenvalues. The derivative of the unitary polar factor solves a Sylvester
equation, which is used in the perturbation theory of the polar
decomposition \cite{bha97}; we are not aware of a statement on the sign of the derivative
of its argument. The phases of a sectorial matrix in the sense of
\cite{wckq20} are defined through a congruence and are a different object.

\begin{corollary}\label{cor:comparison}
Let $\Ima\Omega^{(0)}>0$, let $K$ be Hermitian, and let $\Phi^{(0)}$ and $\Phi$
be the polar phases of $\Omega^{(0)}$ and $\Omega^{(0)}+K$. Then
$\tr(K\Phi)\leq\tr(K\Phi^{(0)})$. In particular, if
$\Omega^{(0)}=zP$ with $\Ima z>0$ and $P>0$, then
$\tr(K\Phi)\leq(\arg z)\tr K$.
\end{corollary}

\begin{proof}
Theorem~\ref{thm:phase} with $\Omega'=\Omega^{(0)}$ and
$\Omega''=\Omega^{(0)}+K$. The polar phase of $zP=(|z|P)e^{\ii\arg z}$ is the
scalar $\arg z$.
\end{proof}

\begin{corollary}\label{cor:blocks}
Let $\CC^n=H\oplus H'$ be an orthogonal decomposition and
\[
 \Omega=\begin{pmatrix}\Omega_{11}&K_{12}\\K_{12}^\dagger&\Omega_{22}\end{pmatrix},
 \qquad\Ima\Omega>0\text{ block diagonal},
\]
with polar phase $\Phi=(\Phi_{ik})$. Then
$\Rea\tr(K_{12}\,\Phi_{21})\leq0$, with equality only if $K_{12}=0$.
\end{corollary}

\begin{proof}
Let $K$ be the Hermitian matrix with off-diagonal blocks $K_{12},K_{12}^\dagger$
and zero diagonal blocks, $\Omega(\varepsilon)=\Omega-(1-\varepsilon)K$ and
$\Sigma=\diag(\id,-\id)$. Then $\Ima\Omega(\varepsilon)=\Ima\Omega$,
$\Sigma\Omega(\varepsilon)\Sigma=\Omega(-\varepsilon)$, and the polar phase is
equivariant under unitary conjugation; so
$f(\varepsilon)=\tr(K\Phi(\varepsilon))$ is odd, $f(0)=0$, and $f'<0$ for
$K\neq0$ by Theorem~\ref{thm:phase}. Finally
$f(1)=2\Rea\tr(K_{12}\Phi_{21})$.
\end{proof}

The rest of this section proves Theorem~\ref{thm:phase}. The pairing
$\tr(\dot\Omega\dot\Phi)$ is expressed through the solution of a Sylvester
equation, and the resulting quadratic form is a sum of forms with matrices
$(\Ima\Omega)\circ\mathsf L^{(a)}$, where the matrices $\mathsf L^{(a)}$ are
positive semidefinite by an integral representation of the kind that
characterizes anti-L\"owner functions \cite{aud11}.

\subsection{The first variation of the phase}
Put
\[
 \tilde P=e^{-\ii\Phi/2}\,|\Omega|\,e^{\ii\Phi/2}>0,\qquad\text{so that}\qquad
 \Omega=e^{\ii\Phi/2}\,\tilde P\,e^{\ii\Phi/2}.
\]
We work in an orthonormal eigenbasis of $\Phi$, $\Phi=\diag(\phi_a)$ with
$\phi_a\in(0,\pi)$. Then
$\Omega_{ab}=e^{\ii(\phi_a+\phi_b)/2}\tilde P_{ab}$, and since $\tilde P$ is
Hermitian,
\begin{equation}\label{eq:reim}
 (\Rea\Omega)_{ab}=\tilde P_{ab}\cos\tfrac{\phi_a+\phi_b}2,\qquad
 (\Ima\Omega)_{ab}=\tilde P_{ab}\,\varsigma_{ab},\qquad
 \varsigma_{ab}:=\sin\tfrac{\phi_a+\phi_b}2>0 .
\end{equation}

\begin{lemma}\label{lem:variation}
Let $\Omega(\varepsilon)$ be a smooth family with $\Omega(0)=\Omega$,
$\Ima\Omega(\varepsilon)>0$, and derivative $\dot\Omega$ at $0$. There is a
Hermitian matrix $\Theta$ with
\begin{equation}\label{eq:sylvester}
 \tilde P\Theta+\Theta\tilde P=\frac{X-X^\dagger}{2\ii},\qquad
 X=e^{-\ii\Phi/2}\,\dot\Omega\,e^{-\ii\Phi/2},
\end{equation}
and, in the eigenbasis of $\Phi$,
\begin{equation}\label{eq:phidot}
 \dot\Phi_{ab}=\frac{\phi_a-\phi_b}{\sin\frac{\phi_a-\phi_b}2}\,\Theta_{ab},
\end{equation}
where the factor is read as $2$ if $\phi_a=\phi_b$. If $\dot\Omega$ is
Hermitian, the right side of \eqref{eq:sylvester} has entries
$-\dot\Omega_{ab}\,\varsigma_{ab}$.
\end{lemma}

\begin{proof}
Let $\mathsf U(\varepsilon)=e^{\ii\Phi(\varepsilon)/2}$, so that
$\Omega=\mathsf U\tilde P\mathsf U$ for every $\varepsilon$. Since
$\mathsf U$ is unitary,
$\dot{\mathsf U}=\ii\Upsilon\mathsf U=\mathsf U\,\ii\Upsilon'$ with $\Upsilon$ Hermitian and
$\Upsilon'=\mathsf U^{-1}\Upsilon\mathsf U$. Hence
$\dot\Omega=\ii\Upsilon\Omega+\mathsf U\dot{\tilde P}\mathsf U+\ii\Omega\Upsilon'$ and
\[
 \mathsf U^{-1}\dot\Omega\,\mathsf U^{-1}
 =\dot{\tilde P}+\ii\bigl(\Upsilon'\tilde P+\tilde P\Upsilon\bigr).
\]
Apply $M\mapsto(M-M^\dagger)/2\ii$. The Hermitian matrix $\dot{\tilde P}$
drops out and $\ii(\Upsilon'\tilde P+\tilde P\Upsilon)$ contributes
$\frac12[(\Upsilon+\Upsilon')\tilde P+\tilde P(\Upsilon+\Upsilon')]$; so \eqref{eq:sylvester}
holds with $\Theta=\frac12(\Upsilon+\Upsilon')$. For an entire function $f$ and
a smooth family of Hermitian matrices, differentiating the Cauchy integral
$f(\Phi)=\frac1{2\pi\ii}\oint f(\zeta)(\zeta-\Phi)^{-1}d\zeta$ gives, in an
eigenbasis of $\Phi$,
$(\frac d{d\varepsilon}f(\Phi))_{ab}=\dot\Phi_{ab}\,(f(\phi_a)-f(\phi_b))/(\phi_a-\phi_b)$,
with the value $f'(\phi_a)$ if $\phi_a=\phi_b$ (the formula of Daleckii and
Krein; see \cite{bha97}). For $f(x)=e^{\ii x/2}$ and
$\gamma=\phi_a-\phi_b$ this gives
$\Upsilon_{ab}=-\ii\dot{\mathsf U}_{ab}e^{-\ii\phi_b/2}
=\dot\Phi_{ab}\,2\gamma^{-1}\sin(\gamma/4)\,e^{\ii\gamma/4}$ and
$\Upsilon'_{ab}=e^{-\ii\gamma/2}\Upsilon_{ab}$, hence
$\Theta_{ab}=\dot\Phi_{ab}\sin(\gamma/2)/\gamma$. For Hermitian
$\dot\Omega$,
$(X-X^\dagger)_{ab}/2\ii=\dot\Omega_{ab}(e^{-\ii(\phi_a+\phi_b)/2}-e^{\ii(\phi_a+\phi_b)/2})/2\ii
=-\dot\Omega_{ab}\,\varsigma_{ab}$.
\end{proof}

The equation \eqref{eq:sylvester} has a unique solution because
$\tilde P>0$.

\begin{lemma}\label{lem:form}
Let
\[
 \varrho_{ab}=\frac{\phi_a-\phi_b}{2(\cos\phi_b-\cos\phi_a)}\quad(\phi_a\neq\phi_b),
 \qquad\varrho_{ab}=\frac1{2\sin\phi_a}\quad(\phi_a=\phi_b),
\]
a positive symmetric function of $(\phi_a,\phi_b)$. For a family with
constant imaginary part, $\tr(\dot\Omega\dot\Phi)=-4\,Q(\Theta)$, where
\[
 Q(\Theta)=\sum_{a,b}\varrho_{ab}\,(\tilde P\Theta+\Theta\tilde P)_{ba}\,\Theta_{ab}
 =\sum_a\inner{\theta^a}{\bigl((\Ima\Omega)\circ\mathsf L^{(a)}\bigr)\theta^a},\qquad
 \mathsf L^{(a)}_{bd}=\frac{\varrho_{ba}+\varrho_{da}}{\varsigma_{bd}},
\]
and $\theta^a$ is the $a$-th column of $\Theta$.
\end{lemma}

\begin{proof}
By Lemma~\ref{lem:variation},
$\dot\Omega_{ab}=-(\tilde P\Theta+\Theta\tilde P)_{ab}/\varsigma_{ab}$ and
$\dot\Phi_{ab}=\gamma\,\Theta_{ab}/\sin(\gamma/2)$ with
$\gamma=\phi_a-\phi_b$, and
$\gamma/(\sin(\gamma/2)\,\varsigma_{ab})=2\gamma/(\cos\phi_b-\cos\phi_a)=4\varrho_{ab}$.
So
$\tr(\dot\Omega\dot\Phi)=\sum_{a,b}\overline{\dot\Omega_{ab}}\dot\Phi_{ab}=-4Q(\Theta)$,
the conjugate of an entry of a Hermitian matrix being the transposed entry.
Expanding,
$Q=\sum_{a,b,c}\varrho_{ab}[\tilde P_{bc}\Theta_{ca}\Theta_{ab}
+\Theta_{bc}\tilde P_{ca}\Theta_{ab}]$; renaming the indices $(a,b,c)$ of the
second sum as $(c,a,b)$ turns it into
$\sum_{a,b,c}\varrho_{ca}\tilde P_{bc}\Theta_{ca}\Theta_{ab}$. Hence
$Q=\sum_{a,b,c}(\varrho_{ab}+\varrho_{ca})\tilde P_{bc}\,\Theta_{ab}\Theta_{ca}$,
and $\tilde P_{bc}=(\Ima\Omega)_{bc}/\varsigma_{bc}$ by \eqref{eq:reim}, while
$\Theta_{ab}\Theta_{ca}=\overline{\theta^a_b}\,\theta^a_c$.
\end{proof}

\subsection{Positivity of the kernels}

\begin{lemma}\label{lem:integral}
For $x,c>0$ let
$g_c(x)=(1+x^2)(\arctan x-\arctan c)/(x^2-c^2)$ for $x\neq c$ and
$g_c(c)=1/(2c)$. Then
\[
 g_c(x)=\frac2\pi\int_0^\infty\frac x{x^2+y^2}\,w_c(y)\,dy,\qquad
 w_c(y)=\frac{|1-y^2|\,k_c(y)}{y^2+c^2}\geq0,
\]
where $k_c(y)=\arctan c$ for $0<y<1$ and $k_c(y)=\frac\pi2-\arctan c$ for
$y>1$.
\end{lemma}

\begin{proof}
For $x\neq c$,
$(1-y^2)/((x^2+y^2)(c^2+y^2))=p/(x^2+y^2)+p'/(c^2+y^2)$ with
$p=(1+x^2)/(c^2-x^2)$ and $p'=(1+c^2)/(x^2-c^2)$. Moreover
$\int_0^1x\,dy/(x^2+y^2)=\frac\pi2-\arctan x$,
$\int_1^\infty x\,dy/(x^2+y^2)=\arctan x$,
$\int_0^1x\,dy/(c^2+y^2)=\frac xc(\frac\pi2-\arctan c)$ and
$\int_1^\infty x\,dy/(c^2+y^2)=\frac xc\arctan c$. With $\beta=\arctan c$ and
$\beta'=\frac\pi2-\beta$,
\begin{align*}
 \int_0^\infty\frac{x\,w_c(y)}{x^2+y^2}\,dy
 &=\beta\bigl[p(\tfrac\pi2-\arctan x)+p'\tfrac xc\beta'\bigr]
 -\beta'\bigl[p\arctan x+p'\tfrac xc\beta\bigr]\\
 &=\tfrac\pi2\,p\,(\arctan c-\arctan x)=\tfrac\pi2\,g_c(x).
\end{align*}
Both sides are continuous in $x$, which gives the case $x=c$.
\end{proof}

\begin{lemma}\label{lem:antiloewner}
For every $c>0$ and all $x_1,\dots,x_n>0$ the matrix
$\bigl[(g_c(x_b)+g_c(x_d))/(x_b+x_d)\bigr]_{b,d}$ is positive semidefinite.
\end{lemma}

\begin{proof}
For $f_y(x)=x/(x^2+y^2)$,
\[
 \frac{f_y(x)+f_y(x')}{x+x'}=\frac{xx'+y^2}{(x^2+y^2)(x'^2+y^2)}
 =f_y(x)f_y(x')+y^2\,\frac1{x^2+y^2}\,\frac1{x'^2+y^2},
\]
a sum of two kernels of the form $\psi(x)\psi(x')$. Integrate against the
positive measure $\frac2\pi w_c(y)\,dy$ of Lemma~\ref{lem:integral}.
\end{proof}

Functions $g$ on $(0,\infty)$ for which all matrices
$[(g(x_b)+g(x_d))/(x_b+x_d)]$ are positive semidefinite are called
anti-L\"owner; Audenaert \cite{aud11} characterizes them by an integral
representation, and Lemma~\ref{lem:integral} puts $g_c$ in that form. The
sufficient condition of Kwong quoted in \cite{aud11}, operator monotonicity
of $g$, does not apply, since $g_c$ is not monotone: for $c=1$ the limits at
$0$ and at $\infty$ are $\frac\pi4$, while the value at $1$ is $\frac12$.

\begin{lemma}\label{lem:antiloewnermatrix}
For every index $a$ the matrix $\mathsf L^{(a)}$ of Lemma~\ref{lem:form} is
positive semidefinite, with positive diagonal.
\end{lemma}

\begin{proof}
Put $x_b=\tan(\phi_b/2)>0$ and $c=x_a$. Then
$\cos\phi_a-\cos\phi_b=2(x_b^2-c^2)/((1+c^2)(1+x_b^2))$ and
$\phi_b-\phi_a=2(\arctan x_b-\arctan c)$, so
$\varrho_{ba}=\frac{1+c^2}2g_c(x_b)$, also when $\phi_b=\phi_a$; and
$\varsigma_{bd}=(x_b+x_d)/(\sqrt{1+x_b^2}\sqrt{1+x_d^2})$. Hence
\[
 \mathsf L^{(a)}_{bd}=\frac{1+c^2}2\,\sqrt{1+x_b^2}\,\sqrt{1+x_d^2}\;
 \frac{g_c(x_b)+g_c(x_d)}{x_b+x_d},
\]
the matrix of Lemma~\ref{lem:antiloewner} conjugated by a positive diagonal
matrix. The diagonal entries $2\varrho_{ba}/\varsigma_{bb}$ are positive.
\end{proof}

\begin{proof}[Proof of Theorem~\ref{thm:phase}]
By Lemma~\ref{lem:form},
$\tr(\dot\Omega\dot\Phi)=-4\sum_a\inner{\theta^a}{((\Ima\Omega)\circ\mathsf L^{(a)})\theta^a}$.
Let $y_{\min}>0$ be the least eigenvalue of $\Ima\Omega$. Since
$\Ima\Omega-y_{\min}\id\geq0$ and $\mathsf L^{(a)}\geq0$, the Schur product
theorem \cite[Section~7.5]{hj13} gives
$(\Ima\Omega)\circ\mathsf L^{(a)}\geq y_{\min}\diag(\mathsf L^{(a)}_{bb})>0$ by
Lemma~\ref{lem:antiloewnermatrix}. So $Q(\Theta)>0$ unless $\Theta=0$, and
$\Theta=0$ forces $\dot\Omega=0$ by \eqref{eq:sylvester}. The integrated
statement follows along the segment
$\Omega'+\varepsilon(\Omega''-\Omega')$, on which the imaginary part is
constant.
\end{proof}

\begin{remark}\label{rem:phasefails}
The positivity of $\Ima\Omega$ is needed. (The numbers of this remark are
computed in floating-point arithmetic and are not used in any proof.) The
polar phase is defined and
analytic near any invertible matrix whose unitary polar factor has spectrum
in the open upper half-circle. For $\Phi=\diag(\frac15,3)$,
$\tilde P=\left(\begin{smallmatrix}1&9/10\\9/10&1\end{smallmatrix}\right)$ and
$\Omega=e^{\ii\Phi/2}\tilde Pe^{\ii\Phi/2}$, the matrix $\Ima\Omega$ has
eigenvalues $-0.730\ldots$ and $1.0699\ldots$, and the symmetric part of the
derivative of the phase on Hermitian variations has the positive eigenvalue
$0.2335\ldots$ For $\Ima\Omega>0$ the phase is in general not the gradient of
a function of $\Rea\Omega$ for the trace inner product: at
$\Omega=\left(\begin{smallmatrix}\ii&1\\1&3\ii\end{smallmatrix}\right)$ the
derivative of the entry $\Phi_{11}$ with respect to $(\Rea\Omega)_{22}$ is
$0.0488\ldots$, while that of $\Phi_{22}$ with respect to $(\Rea\Omega)_{11}$
is $0.0306\ldots$ Nor is $\Rea\Omega\mapsto-\Phi$ monotone for the L\"owner
order, even for a scalar imaginary part $y$: there
$-\Phi=\arctan(\Rea\Omega/y)-\frac\pi2$, and the arctangent is not operator
monotone on the real line, since an operator monotone function on the real
line is affine. The
trace norm $\tr|\Omega|$ is a convex function of $\Rea\Omega$ with gradient
$\cos\Phi$, so $\cos\Phi$ is a monotone map; Theorem~\ref{thm:phase} concerns
$\Phi$ itself and does not follow from this.
\end{remark}
